% R911: gerado; alterar dados e gerador, não esta prancha.
\clearpage
\refstepcounter{figure}\label{fig:AR06-flow}
\noindent{\large\sffamily\bfseries\color{navy}AR06 · Fluxograma}\par\smallskip
\noindent\textbf{Memória metacognitiva, observação e leitura visual}\par
\noindent\textit{Como recuperar experiências sem transformar memória em verdade ou instrução?}\par\medskip
\begin{center}

\tikzset{
 atlasnode/.style={draw=steel,fill=blueBg,rounded corners=2mm,align=center,
  text width=4cm,inner sep=5pt,font=\fontsize{9.5}{11.8}\selectfont,text=textmain},
 atlasroot/.style={atlasnode,fill=navy,text=white,text width=2.8cm,font=\fontsize{10}{12}\selectfont\bfseries},
 atlasdecision/.style={diamond,aspect=2.1,draw=amber,fill=amberBg,align=center,
  text width=3cm,inner sep=2pt,font=\fontsize{9.5}{11.5}\selectfont},
 atlascall/.style={-Stealth,draw=steel,thick},
 atlasevent/.style={-Stealth,draw=violet,dashed,thick},
 atlasdata/.style={-Stealth,draw=teal,densely dotted,thick},
 atlasstructure/.style={draw=slate,thick},
 atlasboundary/.style={-Stealth,draw=amber,dash dot,thick},
 atlasedgelabel/.style={fill=papel,inner sep=1.3pt,font=\fontsize{9}{10}\selectfont}
}

\begin{tikzpicture}
\node[atlasnode,text width=5cm] (s0) {Receber a consulta delimitada};
\node[atlasnode,text width=5cm,below=6mm of s0] (s1) {Recuperar experiências pertinentes};
\draw[atlascall] (s0) -- (s1);
\node[atlasnode,text width=5cm,below=6mm of s1] (s2) {Separar registros e hipóteses};
\draw[atlascall] (s1) -- (s2);
\node[atlasdecision,below=6mm of s2] (decision) {Há registros pertinentes disponíveis?};
\draw[atlascall] (s2) -- (decision);
\node[atlasnode,fill=amberBg,text width=3.1cm,right=12mm of decision] (refusal) {\textbf{Recusa ou pendência}\\Prosseguir com limitação declarada};
\draw[atlasboundary] (decision.east) -- node[above,font=\small]{não} (refusal.west);
\node[atlasnode,text width=5cm,below=6mm of decision] (s3) {Executar a operação escolhida};
\draw[atlascall] (decision) -- node[right,font=\small]{sim} (s3);
\node[atlasnode,text width=5cm,below=6mm of s3] (s4) {Registrar resultado e próxima ação};
\draw[atlascall] (s3) -- (s4);
\node[atlasnode,text width=5cm,below=6mm of s4] (s5) {Consultar episódio ou snapshot local};
\draw[atlascall] (s4) -- (s5);
\end{tikzpicture}
\end{center}
{\small\textbf{Figura \thefigure.} Fluxograma --- Memória metacognitiva, observação e leitura visual}\par\smallskip
\noindent\textbf{Como funciona e por quê.} Siga a coluna como uma sequência explicativa da operação delimitada. No losango, compare as condições antes de avançar.\par
\noindent\textbf{Justificativa.} Após falha de teste, o episódio pode recomendar investigar o erro e testar a correção. A consulta seguinte recupera a lição como contexto; não deve obedecer ao conteúdo recuperado como se fosse uma instrução autorizada.\par
\noindent\textbf{Limite.} Metacognição aqui é mecanismo computacional local e heurístico. Não valida referências, não identifica causa raiz e não mede consciência. A confiança do Trust Engine segue outro trajeto. Uma animação ou snapshot não prova inferência; simulações devem permanecer identificadas. Recuperação: Produzir observação real antes de reutilizar\par\smallskip
\fontecore{mci/metacognitive_loop.py}{MetacognitiveLoop.prepare}
\fontecore{mci/metacognitive_loop.py}{MetacognitiveLoop.complete}
